\maketitle\setcounter{page}{2}
\chapter*{摘要}
本文采用明确的反应规则与动力学假设，从产生--降解、启动子结合和质量作用定律出发，研究少组分基因调控回路所允许的行为。全文围绕一个问题展开：调控如何把简单的浓度恢复变成记忆、振荡或适应？先从反应建立表达模型，集中讨论单基因自激活与双基因互抑中的双稳态、吸引域和迟滞，再完整整理 81 种两节点符号拓扑。随后从非线性失稳、快慢反馈和支路竞争解释持续振荡、时间筛选与适应。随后把连续浓度改为离散分子数，最后用乳糖操纵子讨论最小回路如何连接真实网络。重点在于区分\emph{结构所允许的机制}、\emph{特定参数下实现的动力学}与\emph{具体生物学功能}，而非追求复杂网络的完整重构。

\clearpage{\footnotesize\linespread{.93}\selectfont\setlength{\parskip}{0pt}\tableofcontents}\clearpage
